\chapter{Mouse encoding and SGR mode}
\label{ch:mouse}

\begin{aside}
Mouse support in terminals is a stack of three encodings, each
fixing the previous. \maya{} requests the newest, falls back to
the next, and exposes a clean event type regardless of which
arrived on the wire.
\end{aside}

\section{Three encodings}

\begin{enumerate}[leftmargin=*]
  \item \textbf{X10 encoding (1984).} Three bytes after the
    \code{ESC [ M} prefix: button + 32, column + 32, row + 32.
    Coordinates capped at 223 (255 minus the 32 offset).
  \item \textbf{UTF-8 encoding.} Same as X10 but coordinates
    encoded as UTF-8, so larger terminals work. Hard to parse
    reliably.
  \item \textbf{SGR encoding (1995, widely adopted ~2008).}
    \code{ESC [ < button ; column ; row M} or \code{m} (the
    final letter is \code{M} for press, \code{m} for release).
    Coordinates are decimal text; no limit.
\end{enumerate}

\maya{} requests SGR (\code{ESC [ ? 1006 h}) and falls back to
X10 only if the terminal rejects it. SGR parsing is simple text
parsing; X10 parsing must handle the 223-cap edge.

\section{Modes}

Mouse modes control \emph{which} events the terminal sends:

\begin{itemize}[leftmargin=*]
  \item \code{? 1000 h}: button press and release only.
  \item \code{? 1002 h}: above + drag (motion with button held).
  \item \code{? 1003 h}: all motion, even without buttons.
\end{itemize}

\code{1003} is expensive on the wire --- a moving mouse fires
events as fast as the terminal can read them. Use \code{1002} for
most apps; reserve \code{1003} for hover-driven UIs.

\maya{} defaults to \code{1002} and offers a flag for
\code{1003}.

\section{Decoding}

A press of button 1 at column 10, row 5 in SGR mode arrives as:

\begin{lstlisting}
ESC [ < 0 ; 10 ; 5 M
\end{lstlisting}

The button code is a small integer with bit flags for shift,
alt, and control modifiers:

\begin{itemize}[leftmargin=*]
  \item Bits 0--1: button (00 = left, 01 = middle, 10 = right,
    11 = release).
  \item Bit 2: shift.
  \item Bit 3: alt.
  \item Bit 4: control.
  \item Bit 5: motion.
  \item Bit 6: scroll (wheel up = 64, wheel down = 65, side =
    66, 67).
\end{itemize}

\maya{}'s decoder normalises into a \code{MouseEvent}:

\begin{lstlisting}
struct MouseEvent {
    int x, y;
    enum Kind { Press, Release, Motion, Wheel } kind;
    enum Button { Left, Middle, Right, None } button;
    int wheel_delta;   // +1 up, -1 down, 0 otherwise
    bool shift, alt, ctrl;
};
\end{lstlisting}

\section{Drag tracking}

Drag is a press followed by motion events with the button bit
still set, followed by release. \maya{} maintains a small state
machine to attach the held-down button to motion events.

\section{Coordinate system}

Mouse coordinates are 1-based, like the cursor. \maya{}
converts to 0-based internally and clips to the canvas.

A click outside the canvas (between resize and renderer
catching up) is reported but the dispatch silently drops it.

\section{Hit-testing}

The event arrives with absolute coordinates. The dispatcher
walks the element tree and asks ``which element owns this
column and row?''.

This is a tree walk with bounding-box check at each node. For
deep trees this could be slow, but typical TUIs are shallow (a
few dozen elements visible at once). The walk is
sub-microsecond.

The hit element receives the event; if it does not consume it,
the event bubbles to the parent.

\section{Modifiers}

Shift-click, Ctrl-click, Alt-click come through as a normal
click event with the corresponding modifier bit. Apps can use
them for ``open in new tab'' style actions.

Some terminals intercept Ctrl-click for their own purposes
(open URL, jump to definition). \maya{} cannot prevent this.

\section{Mouse and selection}

The terminal's built-in selection (drag to highlight, copy
with right-click) is suppressed when \maya{} enables mouse
mode. To allow selection, hold Shift in most terminals --- the
modifier reroutes the drag to the terminal instead of the
application.

\begin{tryit}
Run a \maya{} app and try to copy text by drag. Nothing
happens. Hold Shift and drag --- the terminal selection
activates. This is documented in xterm but rarely advertised.
\end{tryit}

\section{Restoring on exit}

On exit, \maya{} disables all mouse modes:

\begin{lstlisting}
ESC [ ? 1006 l
ESC [ ? 1003 l
ESC [ ? 1002 l
ESC [ ? 1000 l
\end{lstlisting}

Forgetting any one leaves the next program in the terminal
seeing mouse events it does not understand.

\section{Touch and trackpad}

Trackpads emit mouse events. Scroll gestures arrive as wheel
events. Pinch and rotate do not have a standard escape encoding
and are not reported.

iTerm2's ``modifiers as keys'' option converts trackpad pinch
into key sequences. \maya{} sees those as keys, not mouse.

\section{Mouse-less fallback}

Every mouse-driven feature in \maya{} has a keyboard
equivalent. The reason: SSH over a flaky link drops mouse
events more readily than keys, and screen readers do not
synthesise mouse.

A button widget is clickable \emph{and} focusable: enter
activates it.

\section{Recap}

\begin{recap}
\begin{itemize}[leftmargin=*]
  \item SGR mode is the modern encoding; request it.
  \item \code{1002} for press+drag; \code{1003} for hover.
  \item Decode bits give modifiers and motion flags.
  \item Hit-test by walking the element tree.
  \item Hold Shift to fall back to terminal selection.
  \item Always restore on exit.
\end{itemize}
\end{recap}

\section*{Workshop}

\begin{enumerate}[leftmargin=*]
  \item Log every mouse event to a side panel. Move and
    click around; watch the events.
  \item Build a draggable widget: pick-up on press, follow
    on drag, drop on release.
  \item Add a hover effect to a button using mode
    \code{1003}. Measure the wire bandwidth as you wiggle
    the mouse.
\end{enumerate}
